% ----------------------------------------------------------------------
% Conclusão: fica fora das partes, depois de \phantompart
% ----------------------------------------------------------------------
\chapter{Conclusão}
\label{cap:conclusao}

\instrucao{De três a seis páginas. Retome a tese, diga como cada objetivo foi cumprido, explicite a contribuição original para a área, as limitações e uma agenda de pesquisa. Não introduza resultados novos.}

O modelo proposto mostrou-se aplicável e discriminante: separou com clareza as três dimensões e apontou a semântica como prioridade. A contribuição original da tese é um instrumento replicável, com indicadores definidos e pesos justificados, que qualquer rede de repositórios pode aplicar periodicamente.

Como limitação, a amostra se restringiu a universidades públicas. Pesquisas futuras podem estender o modelo a repositórios de dados e acompanhar a evolução dos índices ao longo do tempo.
